
\title{Unix/Linux}
\section{Unix/Linux}

\begin{frame}
    \label{unix}
    \begin{block}{\centering\Large Föreläsning \arabic{section} --- Unix/Linux}
        Förberedelse inför laboration 1.

        \begin{itemize}
            \item Operativsystem, Unix historik
            \item Filer och kataloger
            \item Kommandon
            \item Filskydd
            \item Kommandotolk
            \item Processer
        \end{itemize}
    \end{block}
\end{frame}


\begin{frame}[fragile=singleslide]
    \frametitle{Operativsystem}
    \begin{itemize}
        \item Dator --- kör program
        \item Operativsystem --- en samling program som gör det möjligt att köra ''vanliga'' program
        \item Operativsystemet hanterar:
              \begin{itemize}
                  \item de program som körs (program körs ofta parallellt, operativsystemet ser till att programmen får minne och tid att exekvera)
                  \item yttre enheter (tangentbord, mus, skärm, nätverk, \ldots)
                  \item lagring av data (filer, \ldots)
                  \item skydd, felhantering
                        kommunikation med användaren
              \end{itemize}

        \item Exempel:
              \begin{itemize}
                  \item Linux, Windows, Mac OS X, Unix, \ldots
              \end{itemize}
    \end{itemize}
\end{frame}



\begin{frame}[fragile=singleslide]
    \frametitle{Unix historia}
    \begin{itemize}
        \item 1960-talet: Multics --- Multiplexed Information and Computing System --- stort operativsystem för stora datorer. Skulle lösa alla problem.
        \item 1969: Ken Thompson (AT\&T Bell Labs) började skriva Unix, litet operativsystem för små datorer, avsett för programmerare. Idéer från Multics.
        \item 1973: Unix skrevs om i C av Dennis Ritchie, började distribueras.
        \item Senare: mycket utveckling, många varianter (Linux, Mac OS X, NetBSD, A/IX, HP/UX, Irix, Android, IOS\ldots).
        \item Varumärket Unix ägs idag av The Open Group som är ett leverantörs- och teknikoberoende konsortium. Certifierar produkter (Linux ej certifierat).
    \end{itemize}
\end{frame}

\begin{frame}[fragile=singleslide]
    \frametitle{Kännetecken för Unix}
    Grundläggande:
    \begin{itemize}
        \item Kärna --- liten, grundläggande funktioner
        \item Processer, tidsdelning
        \item Filskydd
        \item Fleranvändarsystem
        \item Kommandotolk (terminalfönster)
    \end{itemize}
    Unix är:
    \begin{itemize}
        \item Litet, standardiserat, flyttbart
        \item (Urspringligen) Inte för nybörjare. Enklare nu med grafiska skrivbordsmiljöer.
    \end{itemize}
\end{frame}

\begin{frame}[fragile=singleslide]
    \frametitle{Vilket operativsystem använder du?}

    \begin{figure}
        \centering
        \includegraphics[height=0.5\textheight]{images/Mentimeter_Branding_Logo.png}
    \end{figure}

\end{frame}

\begin{frame}[fragile=singleslide]
    \frametitle{Filer i Unix/Linux}
    I en dator måste man kunna lagra
    \begin{itemize}
        \item program (Word, Excel, Java-/Scalakompilator, spelprogram, \ldots)
        \item data (foton, document, programmeringsuppgifter, \ldots)
    \end{itemize}

    \halfblankline

    Man lagrar program och data i \emph{filer}. En fil finns normalt på skivminne och har ett namn.

    \halfblankline

    Filnamn i Unix/Linux:

    \begin{itemize}
        \item \code{namn.tillägg} (\code{Calc.scala}, \code{Calc.class}, \code{brev.txt}, \code{test1}, \ldots)
        \item Tillägget är bara en extra upplysning, bestämmer inte filtypen!
    \end{itemize}
\end{frame}

\begin{frame}[fragile=singleslide]
    \frametitle{Kataloger}
    Varje fil finns i en katalog. Kataloger kan innehålla underkataloger, så man får en hierarkisk struktur, ett \emph{filträd}:

    \begin{center}
        % \includegraphics[width=\textwidth,height=0.7\textheight,keepaspectratio]{images/file-tree.pdf}
        \resizebox{!}{0.7\textheight}{% \documentclass{standalone}
% \usepackage{forest}

% \begin{document}

\begin{forest}
    for tree={
    font=\ttfamily,
    grow=south, % Tree grows down
    s sep=8mm, % Horizontal space between siblings
    l sep=6mm, % Vertical space between parent and child
    child anchor=north,
    parent anchor=south,
    anchor=north,
    align=center,
    % edge path={
    %   \noexpand\path [draw, \forestoption{edge}]
    %   (!u.parent anchor) -- ([yshift=2mm].child anchor) -| (.child anchor) \forestoption{edge label};
    % },
    edge path={
            \noexpand\path [draw, \forestoption{edge}]
            (!u.parent anchor) |- ([yshift=2mm].child anchor) -| (.child anchor) \forestoption{edge label};
        },
    draw, % Draw the border of nodes
    rounded corners=3mm, % Round corners of the nodes
    fill=white, % Fill color of nodes
    % drop shadow, % Optional: adds a shadow to nodes
    minimum height=6mm, % Minimum height of nodes
    minimum width=10mm, % Minimum width of nodes
    inner xsep=2mm, % X Padding within nodes
    inner ysep=1mm, % Y Padding within nodes
    calign=center
    }
    [/
    [var]
    [usr]
    [h
        [\(\ldots\),draw=none]
        [d23]
        [d24
            [a]
            [b
                [\(\ldots\),draw=none]
                [ma1337no-s]
                [uw3307bq-s
                    [Downloads]
                    [pgk
                        [lab1
                            [kojo.scala,draw=none]
                        ]
                        [lab2]
                        [\(\ldots\),draw=none]
                    ]
                    [dod
                        [lab1
                            [\(\ldots\),draw=none]
                            [example.txt,draw=none]
                            [\(\ldots\),draw=none]
                        ]
                        [lab2]
                        [lab3]
                        [lab4]
                    ]
                    [\(\ldots\),draw=none]
                ]
                [\(\ldots\),draw=none]
            ]
            [c]
            [d]
            [\(\ldots\),draw=none]
        ]
        [\(\ldots\),draw=none]
    ]
    [bin]
    [etc]
    ]
\end{forest}

% \end{document}
}
    \end{center}

\end{frame}

% \begin{frame}[fragile=singleslide]
%     \frametitle{Filnamn}
%     I praktiken är det inte så rörigt som det ser ut på den förra bilden!
%     Ett fullständigt (absolut) filnamn börjar från rotkatalogen och man räknar upp alla katalogerna med \code{/} emellan:

%     \begin{Code}
%         /h/d24/b/uw3307bq-s/dod/lab1/example.txt
%     \end{Code}

%     Men operativsystemet håller reda på en \emph{aktuell katalog}, som ungefär betyder \emph{``var i filträdet befinner jag mig just nu?''}. Istället för att ange det absoluta filnamnet, räcker det att ange vägen med start från den aktuella katalogen -- ett s.k. \emph{relativt} filnamn.

%     \halfblankline

%     Om \code{uw3307bq-s} (användarens \emph{hemkatalog}) är aktuell katalog så kommer man åt filen ovan med namnet:

%     \begin{Code}
%         dod/lab1/example.txt
%     \end{Code}

% \end{frame}

\begin{frame}[fragile=singleslide]
    \frametitle{Filnamn och sökvägar}

    \textbf{Absolut filnamn:}
    \begin{Code}
        /h/d24/b/uw3307bq-s/dod/lab1/example.txt
    \end{Code}

    \textbf{Relativt filnamn:}
    \begin{Code}
        dod/lab1/example.txt
    \end{Code}

    \textbf{Termer:}
    \begin{itemize}
        \item Sökväg (eng. path) -- vägen till en fil
        \item Absolut sökväg -- från rotkatalogen
        \item Relativ sökväg -- från aktuell katalog
        \item Aktuell katalog -- \emph{där du är i filträdet}
    \end{itemize}

\end{frame}


% \begin{frame}[fragile=singleslide]
%     \frametitle{Förkortade filnamn}
%     En del filer har förkortade namn:

%     \blankline
%     \begin{tabular}{cl}
%         \url{~}     & (tilde) hemkatalogen för aktuell användare (Alt Gr + $\sim$ ) \\
%         \url{~user} & hemkatalogen för användaren med användarnamnet user           \\
%         \code{.}    & (punkt) aktuell katalog                                       \\
%         \code{..}   & (punktpunkt) katalogen ovanför aktuell katalog i filträdet    \\
%     \end{tabular}

%     \blankline
%     Antag att \code{\url{~}/dod/lab1} är aktuell katalog. Man kan navigera upp och ned i filträdet enligt följande exempel:

%     \begin{GobbleCode}{6}
%         ~/pgk/lab1/kojo.scala
%         (till hemkatalogen, ner i pfk, ner i lab1)

%         ../../pgk/lab1/kojo.scala
%         (upp ett steg, upp ett steg, ner i pgk, ner i lab1)
%     \end{GobbleCode}

% \end{frame}

\begin{frame}[fragile=singleslide]
    \frametitle{Förkortade filnamn}

    \textbf{Förkortningar:}

    Vissa filnamn har speciella betydelser:

    \blankline
    \begin{tabular}{cl}
        \url{~}     & Hemkatalog (tilde) \\
        \url{~user} & Användarens hemkatalog \\
        \code{.}    & Aktuell katalog (punkt) \\
        \code{..}   & Överordnad katalog (punktpunkt) \\
    \end{tabular}
    \blankline

    \textbf{Exempel:}

    Om \code{\url{~}/dod/lab1} är aktuell katalog:

    \begin{GobbleCode}{6}
        ~/pgk/lab1/kojo.scala
        (till hemkatalogen, ner i pfk, ner i lab1)

        ../../pgk/lab1/kojo.scala
        (upp ett steg, upp ett steg, ner i pgk, ner i lab1)
    \end{GobbleCode}

\end{frame}


% \begin{frame}[fragile=singleslide]
%     \frametitle{Kommandotolk}
%     Man kommunicerar med Unix eller Linux genom att ge kommandon. Ett särskilt program i operativsystemet, kommandotolken, ''shell'', läser in kommandona och utför det som ska göras.

%     \pindent Det finns flera olika kommandotolkar. På Linuxdatorerna i E-huset används kommandotolken \code{bash} (''Bourne-Again SHell''). Andra vanliga kommandotolkar är \code{sh}, \code{zsh} och \code{csh}.

%     \pindent Kommandotolken läser kommandon som ges i ett terminalfönster på skärmen. Så här arbetar tolken:

%     \begin{GobbleCode}{4}
%         Upprepa i all oändlighet:
%             Läs ett kommando
%             Om kommandot är ett riktigt kommando:
%                 utför det som ska göras
%             annars:
%                 skriv ut felmeddelande
%     \end{GobbleCode}

% \end{frame}

\begin{frame}[fragile=singleslide]
    \frametitle{Kommandotolk (shell) och terminal}

    Termerna \emph{kommandotolk} och \emph{terminal} används ofta synonymt, men har olika betydelser:
    \halfblankline

    \begin{description}
        \item [Kommandotolk] För kommunikation med operativsystemet, via kommandon
        \item [Terminal] Program som låter användaren skriva kommandon till kommandotolken
        \item [Shell] Engelska för kommandotolk
    \end{description}

    \blankline
    
    Exempel på shells:
    \begin{itemize}
        \item \code{bash} (Bourne-Again SHell)
        \item \code{sh}, \code{zsh}, \code{csh}
    \end{itemize}

\end{frame}

\begin{frame}[fragile=singleslide]
    \frametitle{Kommandotolk (shell) och terminal}

    Så här fungerar en kommandotolk:
    \begin{GobbleCode}{4}
        Upprepa i all oändlighet:
            Läs ett kommando
            Om kommandot är ett riktigt kommando:
                utför det som ska göras
            annars:
                skriv ut felmeddelande
    \end{GobbleCode}

\end{frame}

\begin{frame}[fragile=singleslide]
    \frametitle{Bash vs. Zsh för Mac-användare}

    \textbf{Bash och Zsh:}
    \begin{itemize}
        \item Bash är standard i de flesta Linux-system.
        \item Zsh är standard i macOS från Catalina.
        \item Båda är kraftfulla skal i Unix-system.
        \item För grundläggande användning, t.ex. denna kurs, fungerar båda bra.
    \end{itemize}

    \textbf{Att byta till Bash (valfritt):}
    \begin{itemize}
        \item Öppna Terminal och skriv \code{bash} för att temporärt byta till Bash.
        \item Installera senaste Bash via Homebrew: \code{brew install bash}
    \end{itemize}

    \textbf{Skillnader:}
    \begin{itemize}
        \item Zsh upplevs av många som mer användarvänligt. Erbjuder fler funktioner och anpassningsmöjligheter, men kan vara överväldigande för nybörjare.
        \item Bash är mer konservativt och har bättre kompatibilitet.
    \end{itemize}

\end{frame}

\begin{frame}[fragile=singleslide]
    \frametitle{Unix-shellalternativ för Windows-användare}

    \textbf{Git Bash:}
    \begin{itemize}
        \item Enkel Bash-tillgång, minimal konfiguration.
        \item Installeras med Git för Windows.
    \end{itemize}

    \textbf{WSL (Windows Subsystem for Linux):}
    \begin{itemize}
        \item Kör fullständig Linux i Windows.
    \end{itemize}

    \textbf{Cygwin:}
    \begin{itemize}
        \item Ger bred Unix-vertygsuppsättning.
        \item Mer konfiguration.
    \end{itemize}

    \blankline
    \textbf{Vilket ska jag välja?}
    \begin{itemize}
        \item \textbf{Git Bash}: Grundläggande Bash och Git. (Rekommenderas)
        \item \textbf{WSL}: Fullständig Linux-utveckling. (Rekommenderas)
        \item \textbf{Cygwin}: Omfattande Unix-funktionalitet. (Avråds)
    \end{itemize}

\end{frame}





\begin{frame}[fragile=singleslide]
    \frametitle{Kommandoformat}
    Varje kommando skrivs på en rad:

    \begin{GobbleCode}{6}
        kommandonamn -option1 -option2 ... argument1 ...
    \end{GobbleCode}

    \begin{itemize}
        \item Kommandot talar om vad som ska göras.
        \item Argumenten är (oftast) filer eller kataloger som påverkas av kommandot.
        \item Optionerna modifierar kommandot på något sätt.
    \end{itemize}

    \halfblankline
    Exempel:

    \halfblankline
    \begin{tabular}{lp{60mm}}
        \code{scalac Calc.scala}  & Kompilerar \code{Calc.scala} med Scalakompilatorn                                         \\
        \code{cp -i report.tex old.tex} & Kopierar \code{report.tex} till \code{old.tex}. \code{-i} betyder att systemet frågar om \code{old.tex} ska skrivas över, om den redan finns. \\
    \end{tabular}
\end{frame}

\begin{frame}[fragile=singleslide]
    \frametitle{Kommandon för filhantering}
    Exempel på kommandon för att hantera filer:

    \blankline
    \begin{tabular}{lp{8.5cm}}
        \code{cp orig kopia} & Kopiera filen \code{orig} till \code{kopia}.                                                                                                                                 \\
        \code{less fil}      & Skriv ut \code{fil} på skärmen, en sida i taget.                                                                                                                             \\
        \code{ls [-la] kat}  & Skriv ut en innehållsförteckning över katalogen \code{kat} (aktuell katalog om ingen katalog ges). \code{-l} skriv i ett längre format, \code{-a} skriv också ut punktfiler. \\
        \code{mv fil1 fil2}  & Döp om \code{fil1} till \code{fil2}. Om \code{fil2} är en katalog så flyttas filen till den katalogen.                                                                       \\
        \code{rm fil1 ...}   & Tag bort de angivna filerna.                                                                                                                                                 \\
    \end{tabular}
\end{frame}

\begin{frame}[fragile=singleslide]
    \frametitle{Kataloghantering}
    Exempel på kommandon för att hantera kataloger:

    \blankline
    \begin{tabular}{lp{8.3cm}}
        \code{cd kat}    & Ändra aktuell katalog till \code{kat}. Utan argument blir det hemkatalogen (samma som \code{cd} \url{~}). \\
        \code{mkdir kat} & Skapa underkatalogen \code{kat} i den aktuella katalogen.                                                 \\
        \code{pwd}       & Skriv ut namnet på aktuell katalog.                                                                       \\
        \code{rmdir kat} & Tag bort \code{kat}. Man måste först ta bort alla filerna i katalogen.                                    \\
    \end{tabular}
\end{frame}

% \begin{frame}[fragile=singleslide]
%     \frametitle{Filskydd 1}
%     Varje fil har en ägare. Ägaren identifieras med användarnamn och grupp, till exempel \code{ma1337no-s} i gruppen \code{students}.

%     \pindent Ägaren kan skydda filer så att andra inte kan komma åt dem, eller göra dem publika så att andra kan komma åt dem.

%     \pindent För att ta reda på filskydd och ägare (och storlek och datum) för filer kan man använda kommandot \code{ls} med optionen \code{l}, alltså \code{ls -l}:
%     \begin{GobbleCode}{6}
%         hacke-3{ma1337no-s}: ls -l
%         -rw-r-----  1 ma1337no-s  students  4940 aug 21 11:14 example.txt
%          (skydd)      (ägare)   (grupp)
%     \end{GobbleCode}

% \end{frame}


\begin{frame}[fragile=singleslide]
    \frametitle{Filskydd 1}

    \textbf{Filägande:}
    \begin{itemize}
        \item Ägare identifieras via användarnamn och grupp (t.ex., \code{ma1337no-s} i \code{students})
    \end{itemize}

    \textbf{Filsäkerhet:}
    \begin{itemize}
        \item Ändra åtkomsträttigheter för att skydda eller dela filer
    \end{itemize}

    \textbf{Visa detaljer:}
    \begin{GobbleCode}{6}
        hacke-3{ma1337no-s}: ls -l
        -rw-r-----  1 ma1337no-s  students  4940 aug 21 11:14 example.txt
         (skydd)       (ägare)     (grupp)
    \end{GobbleCode}

\end{frame}


\begin{frame}[fragile=singleslide]
    \frametitle{Filskydd 2}
    \begin{GobbleCode}{6}
        -rw-r-----  1 ma1337no-s  students  4940 aug 21 11:14 example.txt
          u  g  o
    \end{GobbleCode}

    Tre kategorier av användare:

    \begin{description}
        \item[\code{u}] (user) filens ägare
        \item[\code{g}] (group) medlemmar i samma grupp som ägaren
        \item[\code{o}] (others) alla andra
    \end{description}

    Tre olika rättigheter:

    \begin{description}
        \item[\code{r}] (read) tillstånd att läsa
        \item[\code{w}] (write) tillstånd att skriva
        \item[\code{x}] (execute) tillstånd att exekvera program (för kataloger betyder \code{x} tillstånd att titta på innehållet i katalogen)
    \end{description}

\end{frame}

\begin{frame}[fragile=singleslide]
    \frametitle{Filskydd 3}
    Kommando för att ändra filskydd:

    \begin{GobbleCode}{4}
        chmod skydd fil1 fil2 ...
    \end{GobbleCode}

    Filskydd kan anges symboliskt, till exempel \code{u=rw,o=r}, men också numeriskt, där man anger skydden som är satta med ettor, de som inte är satta med nollor. Till exempel:

    \begin{GobbleCode}{4}
        x = 001, w = 010, r = 100, rx = 101, rw = 110, rwx = 111
    \end{GobbleCode}


    Siffrorna tolkar man sedan som binära tal, vilket ger:

    \begin{GobbleCode}{4}
        x = 1, w = 2, r = 4, rx = 5, rw = 6, rwx = 7
    \end{GobbleCode}


    Fullständiga numeriska filskydd blir till exempel:

    \begin{GobbleCode}{4}
        604 <=> u=rw,o=r    705 <=> u=rwx,o=rx
    \end{GobbleCode}
\end{frame}

\begin{frame}[fragile=singleslide]
    \frametitle{bash-finesser 1}
    Man kan editera den aktuella kommandoraden:

    \blankline
    \begin{tabular}{p{2.2cm}p{8cm}}
        \commandchar{Delete}                      & Ta bort tecknet till vänster om markören.                                 \\
        \commandchar{Control-U}                   & Radera hela raden.                                                        \\
        \commandchar{$\leftarrow$\ $\rightarrow$} & Flytta markören på raden.                                                 \\
        \commandchar{Tab}                         & Filnamnskomplettering (man behöver bara skriva början av ett långt namn). \\
    \end{tabular}

    \blankline
    \code{bash} håller reda på de senaste kommandona som utförts, och man kan få tillbaka gamla kommandon:

    \blankline
    \begin{tabular}{p{2.2cm}p{8cm}}
        \code{$\uparrow$} \code{$\downarrow$} & Bläddra bland tidigare kommandon.                 \\
        \code{!abc}                           & Gör om senaste kommando som inleds med \code{abc}. \\
        \code{!!}                             & Gör om senaste kommando.                          \\
    \end{tabular}
\end{frame}

\begin{frame}[fragile=singleslide]
    \frametitle{bash-finesser 2}
    När man refererar till filer kan man utnyttja jokertecken (wildcards):

    \blankline
    \begin{tabular}{ll}
        \code{?} & motsvarar \emph{ett} godtyckligt tecken \\
        \code{*} & motsvarar 0--flera godtyckliga tecken   \\
    \end{tabular}

    \blankline
    Antag att vi har följande filer i vår katalog:
    \begin{GobbleCode}{4}
        Test1.java  Test2.java  Test23.java  Final.java
        Test1.class Test2.class Final.class  FinalReport.tex
        Outline.tex
    \end{GobbleCode}

    Då kan vi till exempel skriva:

    \begin{GobbleCode}{4}
        javac Test?.java (javac Test1.java Test2.java)
        rm *.class       (rm Test1.class Test2.class Final.class)
        gedit F*.tex     (gedit FinalReport.tex)
    \end{GobbleCode}

\end{frame}

% \begin{frame}[fragile=singleslide]
%     \frametitle{Initieringsfiler}
%     Många program läser en initieringsfil när de startas. Initieringsfilen innehåller inställningar för programmet. Namnen på initieringsfilerna inleds med punkt, så man ser normalt inte dessa filer när man gör \code{ls}.

%     \pindent \code{bash} läser initieringsfilen \code{.bash\_profile} när ett nytt terminalfönster skapas. Bra rader att ha i denna fil (gör att man alltid får frågor om man verkligen vill ta bort/skriva över filer):

%     \begin{Code}
%         alias rm='rm -i'
%         alias cp='cp -i'
%         alias mv='mv -i'
%     \end{Code}

%     Varning: kopiera inte andras initieringsfiler om du inte begriper vad som står i dem!
% \end{frame}

\begin{frame}[fragile=singleslide]
    \frametitle{Initieringsfiler}

    \textbf{Om initieringsfiler:}
    \begin{itemize}
        \item Används ofta för att konfigurera program.
        \item Filnamn börjar med punkt och är dolda, s.k. \emph{punktfiler}.
    \end{itemize}
    
    \textbf{\code{bash} initieringsfil:}
    \begin{itemize}
        \item bash läser \code{.bash\_profile} (från hemkatalogen) när ett nytt terminalfönster öppnas.
        \item Nyttiga alias för att förhindra oavsiktliga filoperationer:
    \end{itemize}

    \begin{Code}
        alias rm='rm -i'
        alias cp='cp -i'
        alias mv='mv -i'
    \end{Code}

    \textbf{Varning:}
    \begin{itemize}
        \item Kopiera inte initieringsfiler utan att förstå innehållet!
    \end{itemize}

\end{frame}


% \begin{frame}[fragile=singleslide]
%     \frametitle{Vad är egentligen ett kommando?}
%     Kommandon är av två slag:

%     \begin{description}
%         \item[''inbyggda''] små kommandon som exekveras direkt i kommandotolken, till exempel \code{cd}, \code{pwd}, \code{exit}
%         \item[''vanliga program''] kommandotolken letar efter ett program med samma namn som kommandot
%     \end{description}

%     Kommandotolken letar efter program i de kataloger som anges i sökvägen (operativsystemvariabeln \code{PATH}, som man själv kan sätta med kommandot \code{export}). När man till exempel skriver \code{javac X.java} exekveras programmet \code{/usr/bin/javac}, eftersom katalogen \code{/usr/bin} finns i sökvägen.

%     \pindent När man ska exekvera ett eget program som finns i den aktuella katalogen skriver man:
%     \begin{Code}
%         ./programnamn
%     \end{Code}

%     \code{./} är nödvändigt eftersom \code{.} (den aktuella katalogen) inte finns i sökvägen.
% \end{frame}

\begin{frame}[fragile=singleslide]
    \frametitle{Vad är egentligen ett kommando?}

    \textbf{Typer av kommandon:}
    \begin{description}
        \item[Inbyggda] Direkt exekverade i kommandotolken (t.ex. \code{cd}, \code{pwd}, \code{exit}).
        \item[Vanliga program] Kommandotolken söker efter programfiler i \code{PATH}.
    \end{description}

    \textbf{Programexekvering:}
    \begin{itemize}
        \item \code{PATH} definierar var program letas upp (\code{export PATH=...} för anpassning).
        \item Ex: \code{/usr/bin/javac} finns i \code{PATH}, så \code{javac X.java} går att köra.
    \end{itemize}

    \textbf{Köra egna program:}
    \begin{Code}
        ./programnamn
    \end{Code}
    Använd \code{./} för program i den aktuella katalogen, eftersom \code{.} inte ingår i \code{PATH}.

\end{frame}


% \begin{frame}[fragile=singleslide]
%     \frametitle{Processer}
%     Varje program som körs i Unix körs som en egen process som exekverar självständigt, ''parallellt'' med andra processer. Varje process körs några millisekunder, sedan får nästa process exekvera, osv. Detta hanteras av operativsystemet.

%     \pindent När man exekverar ett program från ett terminalfönster ''låser'' programmet fönstret tills det avslutas --- så vill man ofta att det ska fungera.

%     \pindent Men man kan också köra program ''i bakgrunden'' dvs frikopplade från terminalfönstret. Det gör man genom att skriva \code{\&} sist på kommandoraden, till exempel:

%     \begin{Code}
%         gedit example.txt &
%     \end{Code}

% \end{frame}
\begin{frame}[fragile=singleslide]
    \frametitle{Processer}

    \textbf{Grundläggande om Processer:}
    \begin{itemize}
        \item Varje program körs som en egen process, oberoende och ''parallellt'' med andra.
        \item Processväxling sker kontinuerligt, hanteras automatiskt av operativsystemet.
    \end{itemize}

    \textbf{Exekvering i Terminal:}
    \begin{itemize}
        \item Standard: Kommandotolken väntar på programmet. Terminalen framstår som upptagen tills programmet avslutas.
        \item Bakgrundskörning: Använd \code{\&} för att köra programmet frikopplat från terminalen.
    \end{itemize}

    \textbf{Exempel:}
    \begin{Code}
        gedit example.txt &
    \end{Code}
    \textit{Öppnar \code{gedit} i bakgrunden, terminalen blir omedelbart tillgänglig för nya kommandon.}

\end{frame}


% \begin{frame}[fragile=singleslide]
%     \frametitle{Fönstersystem, skrivbordsmiljö}
%     I Unix och Linux är fönstersystemet inte inbyggt (Unix utvecklades långt innan det fanns grafiska skärmar).

%     \pindent De flesta Unix- och Linuxsystem som kör på persondatorer använder fönstersystemet X Window System, som oftast kallas bara X. X gör det möjligt att ha fönster på skärmen, att flytta dem, ikonifiera dem, osv.

%     \pindent Exakt hur fönstren ska se ut och hur man arbetar med dem bestäms av \emph{fönsterhanteraren}, som också är ett separat program. Ovanpå fönsterhanteraren finns skrivbordsmiljön.

%     \pindent Det bästa (enda) sättet att lära sig att använda en skrivbordsmiljö är att träna!
% \end{frame}
\begin{frame}[fragile=singleslide]
    \frametitle{Fönstersystem, skrivbordsmiljö}

    \textbf{Fönstersystem:}
    \begin{itemize}
        \item I Unix och Linux är fönstersystemet inte inbyggt (Unix utvecklades långt innan det fanns grafiska skärmar).
        \item Fönstersystemet X Window System (ofta kallat \emph{X}).
    \end{itemize}

    \textbf{Fönsterhanterare och skrivbordsmiljö:}
    \begin{itemize}
        \item Fönsterhanteraren bestämmer utseende och funktion hos fönstren.
        \item Skrivbordsmiljön bygger vidare på fönsterhanteraren för en fullständig användarupplevelse.
    \end{itemize}

    \textbf{Lära sig skrivbordsmiljön:}
    \begin{itemize}
        \item För många av er är Linux nytt. Bästa sättet att bli bekant med en skrivbordsmiljö är att träna!
    \end{itemize}

\end{frame}
